\documentclass[../candidatura.tex]{subfiles}
\graphicspath{{\subfix{../imagenes/}}}
\begin{document}

\section{Reparaciones}

\avisoesbozo

Para reparar un sintetizador, es necesario entenderlo y que estén diseñadas las partes para ser reemplazables.

\subsection{Soquetes}

Una estrategia para lograr esto es la técnica de soquetes.

En vez de soldar un chip directamente, se puede usar un soquete donde se pose el chip, que permita que si el chip ya no funciona, poder reemplazarlo con rapidez, o al revés, si el resto de la placa se destruye, se puede rescatar el milagro que es el chip.

\subsection{Listas de materiales}

Al incluir la lista de materiales en la construcción de sintetizadores, es posible reemplazar o reponer cualquier pieza rota o faltante.


\subsection{De lo general a lo específico, del esquemático a la placa}

Al explicar la idea detrás de un sintetizador, es posible repararlo cuando tenga algún problema.

Los sistemas computacionales pueden quedar obsoletos, puede que un lenguaje caiga en desuso y sea reemplazado por otros, pero si la idea está bien explicada y documentada, se puede portar a otro lenguaje que lo reemplace y así el instrumento se resiste a la obsolescencia.

Los circuitos integrados pueden dejar de producirse, pero se pueden reemplazar por otros con características similares. Muchas veces no es importante un amplificador operacional en particular y se puede reemplazar por otro similar; incluso en un resistor o un capacitor su valor específico no es crítico: al cambiarlo solamente cambiará algún parámetro de su rango de afinación, de su tiempo de carga y descarga al oscilar, pero no afecta mayormente la capacidad de producir ruido.

\newpage

\end{document}
